% Printed 105; principal PDF 112; Google PDF 116.
\par\clearpage

\thispagestyle{empty}
\vspace*{\fill}
\begin{center}
\LARGE
CONFORMAL DOUBLE PROJECTION\\[.5ex]
OF THE SPHEROID ONTO THE SPHERE\\[.5ex]
AND THE PLANE.
\end{center}
\vspace*{\fill}

% Printed 106; principal PDF 113; Google PDF 117.
\par\clearpage

\thispagestyle{empty}
\null

% Printed 107; principal PDF 114; Google PDF 118.
\par\clearpage

\vspace*{.11\textheight}
\begin{center}{\LARGE P\,O\,S\,T\,H\,U\,M\,O\,U\,S\ P\,A\,P\,E\,R\,S.}\end{center}
\vspace{.05\textheight}
\begin{center}
[I.]\par\bigskip
{\Large THE ELLIPTIC SPHEROID\\
{[MAPPED ONTO THE SPHERE].}}\par\bigskip
\rule{.09\textwidth}{.4pt}\par\bigskip
[1.]
\end{center}

The first problem is to map the surface of the ellipsoid onto
the surface of the sphere.

[Let]
\[
\begin{array}{l@{\;\ldotp\;\ldotp\;\ldotp\;\ldotp\;}l}
\varphi & \text{latitude on the spheroid}\\
\psi & \text{corresponding latitude on the sphere}\\
e & \text{eccentricity}\\
a & \text{equatorial radius.}
\end{array}
\]
\noindent [Then the following holds (cf.\ Vol.~IV, Art.~13, p.~207 and following).]
\[
\frac{\mathrm d\psi}{\cos\psi}
=\frac{(1-ee)\,\mathrm d\varphi}{\cos\varphi(1-ee\sin\varphi^2)}
\]
\noindent or
\[
\operatorname{tang}(45^\circ+\tfrac12\psi)
=\operatorname{tang}(45^\circ+\tfrac12\varphi)
\left(\frac{1-e\sin\varphi}{1+e\sin\varphi}\right)^{\frac12 e}\,.\,\text{Const.}
\]

\noindent Every infinitesimal figure on the sphere is thus similar to the
corresponding figure on the spheroid, and equal to it if the radius of the sphere
is set to \(r=\frac{a\cos\varphi}{\cos\psi\surd(1-ee\sin\varphi^2)}\). For \(r\) to attain a maximum at \(\varphi=P\),
one sets the constant multiplier
\(=\bigl(\frac{1+e\sin P}{1-e\sin P}\bigr)^{\frac12 e}\), which also gives \(\psi=P\) on that parallel.
It is precisely there, then, that the figures attain their most perfect equal-
\par\smallskip
\noindent\hfill \(14^{*}\)\hspace{.12\textwidth}\par

% Printed 108; principal PDF 115; Google PDF 119.
\par\clearpage
\noindent\makebox[12mm][l]{\fontsize{13}{16}\selectfont 108}\makebox[\dimexpr\linewidth-24mm\relax][c]{\small\textbf{POSTHUMOUS PAPERS.}}\makebox[12mm][r]{}\par\medskip

\noindent ity if the radius of the sphere is set
\(=\frac a{\surd(1-ee\sin P^2)}\). At other parallels the linear dimensions
on the spheroid are then to those on the sphere as
\[
\frac{\cos\varphi\surd(1-ee\sin P^2)}{\cos\psi\surd(1-ee\sin\varphi^2)}:1=1:n.
\]

\noindent [From this it follows that
\[
\frac{\mathrm d\log n}{\mathrm d\psi}=\Big]
\frac{\mathrm dn}{n\,\mathrm d\psi}=\frac{\sin\varphi-\sin\psi}{\cos\psi}.
\]

\noindent The image on the sphere of a segment of the geodesic is here a
small circle. If we set its radius \(=R\), then
\[
\operatorname{cotang} R=\frac{\sin\varphi-\sin\psi}{\cos\psi}\sin\zeta
\left[=\frac{\mathrm d\log n}{\mathrm d\psi}\sin\zeta\right],
\]
\noindent [where] \(\zeta\) [denotes the] azimuth [of the geodesic].

[The image is] concave toward the south if \(\sin\varphi>\sin\psi\).

\bigskip
\begin{center}[2.]\end{center}
Let
\[
\begin{aligned}
\sin\varphi&=s\\
\sin\psi&=\sigma,
\end{aligned}
\]
\noindent [so that]
\[
\begin{aligned}
\frac{\mathrm ds}{\mathrm d\sigma}
&=\frac{(1-eess)(1-ss)}{(1-ee)(1-\sigma\sigma)}
=\frac{1-(1+ee)ss+ees^4}{(1-ee)(1-\sigma\sigma)}\\
\frac{\mathrm{dd}s}{\mathrm d\sigma^2}
&=\frac{(1-eess)(1-ss)}{(1-ee)^2(1-\sigma\sigma)^2}
\{(2-2ee)\sigma-(2+2ee)s+4ees^3\}
\end{aligned}
\]
\noindent\hfill etc.\hfill\null\par

\noindent For [\(\varphi=\)] \(\psi=P\), these values become:
\[
\begin{aligned}
\frac{\mathrm ds}{\mathrm d\sigma}
&=\frac{1-ee\sin P^2}{1-ee}\\
\frac{\mathrm{dd}s}{\mathrm d\sigma^2}
&=-\frac{4ee}{(1-ee)^2}(1-ee\sin P^2)\sin P\\
\mathopen{\Big[}\frac{\mathrm d^3s}{\mathrm d\sigma^3}
&=-2ee\frac{1-ee\sin P^2}{(1-ee)^3\cos P^2}
\{3-ee-(3+13ee)\sin P^2+14ee\sin P^4\}\Big].
\end{aligned}
\]
\noindent Consequently
\[
\sin\varphi=\sin P+\frac{1-ee\sin P^2}{1-ee}(\sin\psi-\sin P)
-\frac{2ee}{(1-ee)^2}(1-ee\sin P^2)\sin P(\sin\psi-\sin P)^2-\ldots
\]

% Printed 109; principal PDF 116; Google PDF 120.
\par\clearpage
\noindent\makebox[12mm][l]{}\makebox[\dimexpr\linewidth-24mm\relax][c]{\footnotesize\textbf{CONFORMAL PROJECTION OF THE SPHEROID ONTO THE SPHERE.}}\makebox[12mm][r]{\fontsize{13}{16}\selectfont 109}\par\medskip

\noindent [or]
\[
\begin{gathered}
\sin\varphi-\sin\psi=\frac{ee\cos P^2}{1-ee}(\sin\psi-\sin P)
-\frac{2ee}{(1-ee)^2}(1-ee\sin P^2)\sin P(\sin\psi-\sin P)^2\\
-\left[\frac{ee}{3(1-ee)^3}\,.\,\frac{1-ee\sin P^2}{\cos P^2}
\{3-ee-(3+13ee)\sin P^2+14ee\sin P^4\}(\sin\psi-\sin P)^3\right]\ldots.
\end{gathered}
\]
\noindent Furthermore,
\[
\cos\psi^2=\cos P^2-2\sin P(\sin\psi-\sin P)-[(\sin\psi-\sin P)^2].
\]
\noindent Thus
\[
\begin{aligned}
&\frac{\mathrm d\log n}{\mathrm d\sin\psi}
\left[=\frac{\sin\varphi-\sin\psi}{\cos\psi^2}\right]\\
&=\frac{ee}{1-ee}(\sin\psi-\sin P)-\frac{2e^4\sin P}{(1-ee)^2}(\sin\psi-\sin P)^2
-\frac{e^4(5-ee(3+14\sin P^2))}{3(1-ee)^3}(\sin\psi-\sin P)^3\ldots
\end{aligned}
\]
\noindent and
\[
\begin{aligned}
\log n
&=\frac{ee}{2(1-ee)}(\sin\psi-\sin P)^2
-\frac{2e^4\sin P}{3(1-ee)^2}(\sin\psi-\sin P)^3\ldots\\
&=\alpha(\sin\psi-\sin P)^2-\beta(\sin\psi-\sin P)^3\ldots.
\end{aligned}
\]
\noindent [For the flattening \(\frac1{303{,}45..}\) and \(P=51^\circ31'48''{,}70\), we have]
\[
\begin{array}{@{}r@{\quad}l@{}}
\makebox[25mm][s]{\(\displaystyle\frac{ee}{1-ee}\)\hfill .\hfill .\hfill .\hfill .\hfill .}
&7{,}821\,0885\\[.7ex]
\makebox[25mm][s]{\(\displaystyle\frac{\mathrm{Mod.}}2\)\hfill .\hfill .\hfill .\hfill .\hfill .\hfill .}
&9{,}336\,7543\\[.7ex]\hline
\log\alpha=&7{,}157\,8428\\[.5ex]
\makebox[25mm][s]{\(\displaystyle\frac{ee}{1-ee}\)\hfill .\hfill .\hfill .\hfill .\hfill .}
&7{,}821\,0885\\[.7ex]
\makebox[25mm][s]{\(\tfrac43\)\hfill .\hfill .\hfill .\hfill .\hfill .\hfill .\hfill .\hfill .}
&0{,}124\,9387\\
\makebox[25mm][s]{\(\sin P\)\hfill .\hfill .\hfill .\hfill .\hfill .}
&9{,}893\,7263\\\hline
\log\beta=&4{,}997\,5963
\end{array}
\]
\noindent [so, for Briggsian logarithms,]
\[
\log n=0{,}001\,4382.78(\sin\psi-\sin P)^2-0{,}000\,0099.45(\sin\psi-\sin P)^3\ldots.
\]

\begin{center}[3.]\end{center}
Furthermore,
\[
\begin{aligned}
\frac{\mathrm ds}{\mathrm d\psi}
&=\frac{1-(1+ee)ss+ees^4}{(1-ee)\surd(1-\sigma\sigma)},\\
\frac{\mathrm{dd}s}{\mathrm d\psi^2}
&=\frac{(1-ss)(1-eess)}{(1-ee)^2(1-\sigma\sigma)}
\{(1-ee)\sigma-2(1+ee)s+4ees^3\}
\end{aligned}
\]
\noindent\hfill etc.\hfill\null\par

% Printed 110; principal PDF 117; Google PDF 121.
\par\clearpage
\noindent\makebox[12mm][l]{\fontsize{13}{16}\selectfont 110}\makebox[\dimexpr\linewidth-24mm\relax][c]{\small\textbf{POSTHUMOUS PAPERS.}}\makebox[12mm][r]{}\par\medskip

\noindent Thus for [\(\varphi=\)] \(\psi=P\):
\[
\begin{aligned}
\frac{\mathrm ds}{\mathrm d\psi}
&=\frac{(1-ee\sin P^2)\cos P}{1-ee}\\
\frac{\mathrm{dd}s}{\mathrm d\psi^2}
&=-\frac{(1-ee\sin P^2)\sin P}{(1-ee)^2}(1+3ee-4ee\sin P^2)\\
\mathopen{\Big[}\frac{\mathrm d^3s}{\mathrm d\psi^3}
&=-\frac{(1-ee\sin P^2)\cos P}{(1-ee)^3}
\{1+4ee-e^4-(18ee+14e^4)\sin P^2+28e^4\sin P^4\},
\end{aligned}
\]
\noindent consequently]
\[
\begin{gathered}
\sin\varphi=\sin P+\frac{(1-ee\sin P^2)\cos P}{1-ee}(\psi-P)
-\frac{(1-ee\sin P^2)\sin P}{2(1-ee)^2}(1+3ee-4ee\sin P^2)(\psi-P)^2\\
-\left[\frac{(1-ee\sin P^2)\cos P}{6(1-ee)^3}
\{1+4ee-e^4-(18ee+14e^4)\sin P^2+28e^4\sin P^4\}(\psi-P)^3\right]\ldots.
\end{gathered}
\]
\noindent [Since]
\[
\sin\psi=\sin P+\cos P\,.\,(\psi-P)-\tfrac12\sin P\,.\,(\psi-P)^2
-[\tfrac16\cos P\,.\,(\psi-P)^3]\ldots
\]
\noindent [holds, it follows that]
\[
\begin{aligned}
\sin\varphi-\sin\psi&=\frac{ee\cos P^3}{1-ee}(\psi-P)\\
&\quad-\frac{ee}{2(1-ee)^2}\sin P\{5-ee-(5+3ee)\sin P^2+4ee\sin P^4\}(\psi-P)^2\\
&\quad-\Big[\frac{ee}{6(1-ee)^3}\cos P\{7-4ee+e^4-(19+18ee-e^4)\sin P^2\\
&\qquad\qquad\qquad+(46ee+14e^4)\sin P^4-28e^4\sin P^6\}(\psi-P)^3\Big]\ldots.
\end{aligned}
\]
\noindent [Using]
\[
\cos\psi=\cos P-\sin P\,.\,(\psi-P)-[\tfrac12\cos P\,.\,(\psi-P)^2]\ldots
\]
\noindent [we therefore obtain]
\[
\begin{aligned}
&\frac{\mathrm d\log n}{\mathrm d\psi}
\left[=\frac{\sin\varphi-\sin\psi}{\cos\psi}\right]\\
&\quad=\frac{ee}{1-ee}\cos P^2\,.\,(\psi-P)
-\tfrac12\frac{ee}{1-ee}\sin P\cos P
\left(3+4\frac{ee}{1-ee}\cos P^2\right)(\psi-P)^2\\
&\qquad-\Big\{\frac{ee}{1-ee}(\tfrac23\cos P^2-\tfrac12\sin P^2)
-\left(\frac{ee}{1-ee}\right)^2(4\sin P^2-\tfrac53\cos P^2)\cos P^2\\
&\qquad\qquad\qquad-\left(\frac{ee}{1-ee}\right)^3
(4\sin P^2-\tfrac23\cos P^2)\cos P^4\Big\}(\psi-P)^3\ldots.
\end{aligned}
\]

% Printed 111; principal PDF 118; Google PDF 122.
\par\clearpage
\noindent\makebox[12mm][l]{}\makebox[\dimexpr\linewidth-24mm\relax][c]{\footnotesize\textbf{CONFORMAL PROJECTION OF THE SPHEROID ONTO THE SPHERE.}}\makebox[12mm][r]{\fontsize{13}{16}\selectfont 111}\par\medskip

\begin{center}[4.]\end{center}
[Set
\[
\psi-P=\frac{x\surd(1-ee\sin P^2)}a
\]
\noindent and]
\[
\frac{ee}{1-ee}=h;
\]
\noindent \(x\) [is] positive southward.\quad [Then]
\[
\begin{aligned}
\log n&=\tfrac12h\cos P^2\,\frac{xx(1-ee\sin P^2)}{aa}\\
&\quad+\tfrac12h\sin P\cos P(1+\tfrac43h\cos P^2)
\frac{x^3(1-ee\sin P^2)^{\frac32}}{a^3}\\
&\quad-\{h(\tfrac16\cos P^2-\tfrac18\sin P^2)
-hh(\sin P^2-\tfrac5{12}\cos P^2)\cos P^2\\
&\qquad\qquad-h^3(\sin P^2-\tfrac16\cos P^2)\cos P^4\}
\frac{x^4(1-ee\sin P^2)^2}{a^4}
\end{aligned}
\]
\noindent\hfill etc.\hfill\null\par

\bigskip
\begin{center}[5.]\end{center}
[If \(\xi\) is the value on the spheroid corresponding to \(x\), then
\[
\frac{\mathrm d\xi}{\mathrm dx}=\frac1n=1-\log n+\tfrac12(\log n)^2\ldots,
\]
\noindent so that if]
\[
r=\frac a{\surd(1-ee\sin P^2)}
\]
\noindent [is also set, we obtain:]
\[
\xi=x-\alpha x^3-\beta x^4+\gamma x^5\ldots;
\]
\[
\begin{aligned}
\alpha&=\frac{h\cos P^2}{6rr}\\
\beta&=\frac{h\sin P\cos P}{8r^3}(1+\tfrac43h\cos P^2)\\
\gamma&=\frac1{r^4}\{h(\tfrac1{30}\cos P^2-\tfrac1{40}\sin P^2)
-hh(\tfrac15\sin P^2-\tfrac{13}{120}\cos P^2)\cos P^2\\
&\qquad\qquad-h^3(\tfrac15\sin P^2-\tfrac1{30}\cos P^2)\cos P^4\}.
\end{aligned}
\]
\par
\endgroup
